\section{Conventions and classical estimates}\label{sec:preliminaries}

Throughout the proof,
\[
 x\longrightarrow\infty,\qquad L=\log x,\qquad T=\log L,
 \qquad W=\exp(\sqrt L),\qquad \e(t)=\exp(2\pi i t).
\]
A logarithmic power always has a fixed exponent, chosen before $x$
tends to infinity. The implied constants may depend on all previously
fixed parameters. No effective threshold for $x$ is asserted.
The letters $p$ and $q$ in explicitly indicated prime sums range over
primes; general factorization variables are positive integers.
We put $P^-(1)=\infty$, where $P^-(n)$ is the least prime factor of
$n>1$, and write $\tau(n)$ for its divisor function. The notation
$n\asymp Y$ restricts $n$ to a fixed constant enlargement of a dyad
$[Y,2Y)$. All such enlargements remain bounded independently of $x$.
Dirichlet characters are extended by zero on nonunits.

We use the following classical estimates in the forms recorded in
S~\cite[Section~2]{SmoothShifted}. This explicit list fixes their
uniformity; it does not require any distribution theorem beyond the
stated ranges.

\begin{proposition}[Prime estimates]\label{pre:primes}
The prime number theorem holds with an error smaller than every fixed
negative logarithmic power. Mertens' estimates give
\[
 \sum_{p\le y}\frac1p=\log\log y+O(1),\qquad
 \prod_{p\le y}\left(1-\frac1p\right)
 \sim\frac{e^{-\gamma_E}}{\log y},
\]
where $\gamma_E$ is Euler's constant. For fixed $A,C>0$, uniformly
for $r\le(\log y)^C$ and $(a,r)=1$, the Siegel--Walfisz estimate is
\[
 \#\{p\le y:p\equiv a\pmod r\}
 =\frac{\Li(y)}{\varphi(r)}
   +O_{A,C}\left(\frac{y}{(\log y)^A}\right).
\]
Its constants need not be effective.
\end{proposition}

All short intervals used with the prime number theorem have either a
fixed logarithmic-length interpretation in a fixed positive-power
band, or relative length a fixed negative power of $L$; the absolute
error above is made sufficiently small before taking their
differences. We never assume a relative prime asymptotic on every
arbitrarily short interval.

\begin{proposition}[Large sieve and moment bounds]\label{pre:moments}
For arbitrary coefficients supported on the integers of an interval
of length $H$,
\[
 \sum_{r\le R}\frac r{\varphi(r)}
 \sum_{\chi\bmod r}^{*}
 \left|\sum_n a_n\chi(n)\right|^2
 \ll (H+1+R^2)\sum_n|a_n|^2.
\]
The asterisk denotes primitive characters. For every fixed integer
$k\ge0$ there is $C_k$ such that
\[
 \sum_{n\le Y}\tau(n)^k\ll_k Y(\log(2Y))^{C_k},\qquad
 \sum_{n\le Y}\frac{\tau(n)^k}{n}
 \ll_k(\log(2Y))^{C_k}.
\]
If $P(t)=\sum_{n\le Y}c_n n^{it}$, then on any real interval $I$
of length $H$,
\[
 \int_I|P(t)|^2\,dt
 \ll (H+Y\log(2Y))\sum_n|c_n|^2.
\]
For a set $\mathcal T\subset I$ with mutual spacing at least one,
\[
 \sum_{t\in\mathcal T}|P(t)|^2
 \ll (H+1+Y\log(2Y))(\log(2Y))^2\sum_n|c_n|^2.
\]
The constants in the last two estimates are absolute, independently
of how the coefficients were formed.
\end{proposition}

These are Theorem~2.3 and Lemmas~2.5--2.6 of S. In particular, the
fixed-moment estimate applies to a convolution of a fixed number of
logarithmically bounded sequences, giving coefficient-square mass
$UL^{O(1)}$ on a dyad of size $U$. With reciprocal-index
normalization the bound is $L^{O(1)}/U$. For the one growing power
used later, we instead prove a factorial coefficient bound and use
the coefficient-independent mean-square inequality.

\begin{proposition}[Derivative tests]\label{pre:derivatives}
Let $f$ be real-valued on an interval containing $H$ consecutive
integers. If $f'$ is monotone and remains in
$[j+\lambda,j+1-\lambda]$ for an integer $j$ and
$0<\lambda\le1/2$, then
\[
 \sum_n\e(f(n))\ll\lambda^{-1}.
\]
If instead $f$ is twice continuously differentiable and
$\lambda\le|f''|\le C\lambda$ throughout the interval, then
\[
 \sum_n\e(f(n))\ll_C H\sqrt\lambda+\lambda^{-1/2}.
\]
Both estimates apply on every subinterval satisfying their hypotheses.
\end{proposition}

This is the form of S, Lemma~2.7. The distance from integers in the
first test is part of the hypothesis, not merely a lower bound on
the absolute derivative.

We will repeatedly separate smooth functions in logarithmic
coordinates. The following elementary version of S, Lemma~2.8,
explains the order in which frequency cutoffs and accuracies may be
chosen.

\begin{lemma}[Smooth separation]\label{pre:separation}
Let $F$ be smooth and supported on a fixed bounded box in
$\R^b$, with $b$ fixed. Suppose all derivatives of order $j$ are
bounded by $C_jL^{C(j+1)}$. Fourier inversion separates its variables
with integrated absolute coefficient mass $L^{O(1)}$. There is a
fixed exponent $C'>0$, depending on $b,C$, such that truncating each
frequency at $L^{C'}$ makes an error $O_A(L^{-A})$ for every fixed
$A>0$. The exponent $C'$ is fixed before $A$.
\end{lemma}

\begin{proof}
Integration by parts gives, for each fixed $j$, an integrable tail
bounded by a constant times
$L^{C(j+1)}(1+|\xi|)^{-j}$. Using a fixed $j>b$ first bounds the
full integral by a fixed logarithmic power. Outside
$|\xi|\ge L^{C'}$ the same calculation gives
\[
 O_j\bigl(L^{C(j+1)-C'(j-b)}\bigr).
\]
Fix $C'>C$ and then increase $j$ for any prescribed $A$.
The Fourier phase factors into one-dimensional phases. Applied to
normalized logarithmic coordinates, these are multiplicative
power twists.
\end{proof}

The fixed-dimension condition in Lemma~\ref{pre:separation} will not
be used with a growing-dimensional box. In the allocation argument
there are $O(T)$ separate one-dimensional transitions, each with a
fixed Fourier integral norm. Their product costs only
$\exp(O(T))=L^{O(1)}$; a telescoping estimate controls their combined
tails.

The weighted block sieve and local-density proxies are stated in the
prime-extraction section, where their particular coefficient bounds
and precision dependence are needed. The generic dilation-graph
results are likewise stated at their application in the shifted
correlation section. These are the only nonclassical imported
ingredients of the proof.
